\documentclass[10pt,a4paper]{amsart}
\usepackage[margin=19mm]{geometry}
\usepackage{amsmath,amssymb,mathtools}
\usepackage[scr=rsfs,cal=euler]{mathalfa}
\usepackage{xfrac}
\usepackage[unicode,hidelinks,bookmarks=false]{hyperref}
\newcommand{\ie}{i.e.\ }
\newcommand{\cf}{cf.\ }
\newcommand{\resp}{respectively\ }
\newcommand{\Subsection}{\subsection}
\providecommand{\nobreakdash}{\nobreak\hbox{-}}
\setlength{\emergencystretch}{2em}
\multlinegap0pt
\usepackage[shortlabels]{enumitem}
\usepackage{amssymb}
\usepackage{xcolor}
\usepackage{tikz}
\newtheorem{theorem}{Theorem}
\newtheorem{proposition}{Proposition}
\DeclareMathOperator{\Hom}{Hom}
\newcommand{\evencl}{\mathrm{Cl}_0}
\newcommand{\wmax}{W_{\text{max}}}
\title{\MakeUppercase{Categorical resolutions of $A_2$ singularities}}
\author{C\'eline Fietz}
\date{Source: arXiv:2411.19380v3 (CC BY 4.0)}
\begin{document}
\maketitle

\noindent\emph{Source and licence.} \MakeUppercase{Categorical resolutions of $A_2$ singularities}. Version arXiv:2411.19380v3 (CC BY 4.0), distributed by arXiv under the Creative Commons Attribution 4.0 International licence, http://creativecommons.org/licenses/by/4.0/, as recorded in the arXiv licence metadata for this exact version and shown on the abstract page https://arxiv.org/abs/2411.19380v3. Full licence notice: \texttt{licenses/arxiv-2411.19380-CC-BY-4.0.txt}.\par\smallskip

\noindent\textit{Editorial scope.} Counted unit: the article's proposition technical\_heart (source0.tex lines 988--1006). The article's complete original proof of it is delivered with it (source0.tex lines 1007--1064, one \texttt{proof} environment of 58 lines). No mathematics is written, repaired or re-derived: the statement and the proof are the article's own lines, in the article's own notation, and no retained line is substituted or re-flowed.  Retained uncounted as the source-local context the proof needs: lines 844--846; lines 907--933; lines 929--931.  Nothing was omitted from any retained span, so the package carries no omission record.  No retained line contains a citation, so the wrapper carries no bibliography and no bibliography entry.  No retained line contains a figure environment, an image-inclusion command or drawing code, so no image is delivered and the package declares no asset dependency.  Editorial formatting only, in addition: an AMS \texttt{amsart} preamble that restates the article's own declarations and macros used by the retained lines, namely the environments \texttt{theorem}, \texttt{proposition} on separate counters, together with the standard packages those lines need.  The retained environments print in this document as Theorem 1, Proposition 1 in order of appearance, whereas the article prints the same items with the numbering of its own full document.  The complete native source of the article as distributed in the arXiv e-print for this exact version is bundled unchanged as \texttt{sources/169314/source0.tex}.
\begin{theorem}\label{simples_degenerate}
    Let $(V,q)$ be a quadratic space with $q\neq 0$ and let $K \subset V$ be the kernel of the quadratic form $q$. Then the simple $\evencl(q)$-modules are $I_0^{\wmax}$ and $ I_1^{\wmax}$, up to isomorphism, if $\dim(V/K)$ is even and $I_0^{\wmax} \cong I_1^{\wmax}$ if $\dim(V/K)$ is odd. 
\end{theorem}

\begin{theorem}\label{extensions_and_res}
    Let $(V,q)$ be a quadratic space with $q\neq 0$ of corank $1$ and denote by $\{v_1, \dots, v_{N-1}, \varepsilon\}$ an orthogonal $k$-basis for $V$ such that $v_i^2=1$ for all $i \in \{1,\dots, N-1\}$ and $\varepsilon^2 =0$.
    \begin{enumerate}[label=\roman*)]
        \item Let $N$ be odd and let $S_1,S_2$ be the two simple $\evencl(q)$-modules which are unique up to isomorphism by Theorem \ref{simples_degenerate}. Then there exist non-split extensions of left $\evencl(q)$-modules
        \begin{equation}\label{extensions_comp}
            0 \longrightarrow S_2 \overset{\cdot\varepsilon}{\longrightarrow} P_1 \longrightarrow S_1 \longrightarrow 0
            \quad \text{and} \quad 
            0 \longrightarrow S_1 \overset{\cdot\varepsilon}{\longrightarrow} P_2 \longrightarrow S_2 \longrightarrow 0,
        \end{equation}
        where $P_1,P_2$ are the unique indecomposable projective $\evencl(q)$-modules and the exact sequences are induced by right multiplication with $\varepsilon$. In particular, there exist $2$-periodic projective resolutions 
        \begin{equation}\label{projective_resolutions}
        \begin{aligned}
            \dots \overset{\cdot\varepsilon}{\longrightarrow} P_2 \overset{\cdot\varepsilon}{\longrightarrow} P_1 \overset{\cdot\varepsilon}{\longrightarrow}P_2 \overset{\cdot\varepsilon}{\longrightarrow} P_1 \longrightarrow S_1;\\
            \dots \overset{\cdot\varepsilon}{\longrightarrow} P_1\overset{\cdot\varepsilon}{\longrightarrow} P_2 \overset{\cdot\varepsilon}{\longrightarrow} P_1 \overset{\cdot\varepsilon}{\longrightarrow} P_2 \longrightarrow S_2.
        \end{aligned}
        \end{equation} 
        \item Let $N$ be even and let $S$ denote the simple $\evencl(q)$-module which is unique up to isomorphism by Theorem \ref{simples_degenerate}. Set $j := v_{N-1}\varepsilon$.
        Then there exists a non-split extension of left $\evencl(q)$-modules 
        \begin{equation}
            0 \longrightarrow S \overset{\cdot j}{\longrightarrow} P \longrightarrow S\longrightarrow 0,
        \end{equation}
        where $P := \evencl(q)$ is the unique indecomposable projective $\evencl(q)$-module, up to isomorphism and the exact sequence is induced by right multiplication with $j$. In particular, there exists a $1$-periodic projective resolution
        \begin{equation}\label{even_proj_res}
            \dots \overset{\cdot j}{\longrightarrow} P \overset{\cdot j}{\longrightarrow} P \overset{\cdot j}{\longrightarrow} P \overset{\cdot j}{\longrightarrow} P \longrightarrow S.
        \end{equation}
        \end{enumerate}
\end{theorem}

        \begin{equation}
            \dots \overset{\cdot j}{\longrightarrow} P \overset{\cdot j}{\longrightarrow} P \overset{\cdot j}{\longrightarrow} P \overset{\cdot j}{\longrightarrow} P \longrightarrow S.
        \end{equation}

\begin{proposition}\label{technical_heart}
   In the setting of Theorem \ref{extensions_and_res} the following statements hold.
    \begin{enumerate}[label=\roman*)]
        \item Let $N$ be odd dimensional. Then there exist isomorphisms of $k$-algebras
        \begin{equation*}
        \begin{aligned}
            \mathrm{Ext}^{\bullet}(S_1,S_1) \cong \mathrm{Ext}^{\bullet}(S_2,S_2)  \cong k[\theta],
        \end{aligned}
    \end{equation*}
    where $\theta$ denotes an element of degree $2$.
    \item Let $N$ be even dimensional. Then there exists an isomorphism of $k$-algebras
        \begin{equation*}
        \begin{aligned}
            \mathrm{Ext}^{\bullet}(S,S) \cong k[\theta'],
        \end{aligned}
    \end{equation*}
    where $\theta'$ denotes an element of degree $1$.
    \end{enumerate}
\end{proposition}

\begin{proof}
    Let $N$ be odd dimensional. Considering Theorem \ref{extensions_and_res}(i) we have 
    \begin{equation*}
        \Hom(P_2, S_1) \cong 0 \cong \Hom(P_1, S_2),
    \end{equation*}
    since the extensions (\ref{extensions_comp}) are nontrivial. Moreover, we have isomorphisms
    \begin{equation*}
         \Hom(P_1, S_1) \cong k \cong \Hom(P_2, S_2),
    \end{equation*}
    that follow from the fact that $S_1, S_2$ are simple. More precisely, since $\Hom(S_2,S_1) =0$ and $S_2 \overset{\cdot \varepsilon}{\hookrightarrow} P_1$, we see that $\Hom(P_1, S_1) \cong \Hom(S_1, S_1) \cong k$, because every morphism $P_1 \to S_1$ uniquely factors through a morphism $P_1/S_2\cong S_1 \to S_1$.
    
    We apply the functor $\Hom(-, S_1)$ and $\Hom(-,S_2)$ to the projective resolutions (\ref{projective_resolutions}), respectively. Then, taking  cohomology gives rise to isomorphisms of graded $k$-vector spaces 
    \begin{equation*}
        \mathrm{Ext}^{\bullet}(S_1, S_1) \cong k[\theta] \cong \mathrm{Ext}^{\bullet}(S_2, S_2),
    \end{equation*}
    where $\theta$ denotes a generator of degree $2$.
    To prove that they are in fact isomorphisms of $k$-algebras, we consider the exact triangles
    \begin{equation}\label{sequences_algebra}
        \begin{aligned}
        S_2 \longrightarrow P_1 \longrightarrow S_1 \overset{\eta}{\longrightarrow} S_2[1]
        \quad \text{and} \quad 
        S_1 \longrightarrow P_2 \longrightarrow S_2 \overset{\eta'}{\longrightarrow} S_1[1],
        \end{aligned}
    \end{equation}
    where $\eta, \eta'$ denote nontrivial morphisms of degree $1$ corresponding to the nontrivial extensions $P_1$ and $P_2$. This gives rise to triangles of graded $k[\theta]$-modules
    \begin{equation*}
    \begin{aligned}
        \mathrm{Ext}^{\bullet}(S_2,S_1)[-1] \overset{\eta}{\longrightarrow} \mathrm{Ext}^{\bullet}(S_1,S_1) \longrightarrow k; \\
        \mathrm{Ext}^{\bullet}(S_1,S_1)[-1] \overset{\eta'}{\longrightarrow} \mathrm{Ext}^{\bullet}(S_2,S_1) \longrightarrow k.
    \end{aligned}
    \end{equation*}
    In particular, the maps induced by $\eta$ and $\eta'$ are injective. Therefore, composition with the element $\eta'[1]\circ\eta \in \mathrm{Ext}^{2}(S_1,S_1)$ yields an injective morphism 
    \begin{equation*}
        \mathrm{Ext}^{\bullet}(S_1,S_1)[-2] \longrightarrow \mathrm{Ext}^{\bullet}(S_1,S_1).
    \end{equation*}
    We obtain an isomorphism of $k$-algebras
    \begin{equation*}
        k[\theta] \overset{\cong}{\longrightarrow} \mathrm{Ext}^{\bullet}(S_1,S_1), \quad \theta \mapsto (\eta'[1] \circ\eta\colon S_1 \to S_1[2]).
    \end{equation*}
    To compute $\mathrm{Ext}^{\bullet}(S_2,S_2)$, one applies $\Hom(-,S_2)$ to the exact triangles (\ref{sequences_algebra}) and the statement follows by the same arguments. \\
    Let $N$ be even dimensional and consider the projective resolution (\ref{even_proj_res}). 
    Applying the functor $\Hom(-,S)$ and taking cohomology yields an isomorphism of $k$-vector spaces
    \begin{equation*}
        \mathrm{Ext}^{\bullet}(S,S) \cong k[\theta'], 
    \end{equation*}
    where $\theta'$ denotes a generator of degree $1$. To see that this is in fact an isomorphism of $k$-algebras, we consider the exact triangle 
    \begin{equation*}
         S \longrightarrow \evencl(q) \longrightarrow S\overset{\nu}{\longrightarrow} S[1],
    \end{equation*}
    where $0\neq \nu \in \mathrm{Ext}^1(S,S)$ corresponds to the nontrivial extension $P$. This gives rise to an exact triangle
    \begin{equation*}
        \mathrm{Ext}^{\bullet}(S,S)[-1] \overset{\nu}{\longrightarrow} \mathrm{Ext}^{\bullet}(S,S) \longrightarrow k
    \end{equation*}
    which induces an isomorphism of $k$-algebras 
    \begin{equation*}
        k[\theta'] \overset{\cong}{\longrightarrow}\mathrm{Ext}^{\bullet}(S,S), \quad \theta' \mapsto (\nu\colon S \to S[1]).
    \end{equation*}
\end{proof}

\end{document}
